\chapter{语言、字符、字体与字形}\label{languages}

\section{引言}

\topicindex {languages}

\LUATEX\ 内部对待最终被排版出来的字符与字形的方式，与 \TEX82\ 处理同一批对象的方式
大不相同。解释这种差异最简单的办法，是先聚焦于不受限的水平模式（即段落）和断字。
之后再处理水平模式和数学模式中出现的差异就容易了。

在 \TEX82\ 中，你键入的字符在被主控制循环遇到时会被转换为 \type {char} 节点记录。
\TEX\ 在创建这些记录的同时附加并处理字体信息，因此得到的 ``水平列表
中已经包含连字和隐含字距的最终形态。这种打包是必要的，因为我们可能需要获取
比如一个水平盒子的有效宽度。

当需要对段落中的单词进行断字时，\TEX\ 会（一次处理一个词）把 \type {char} 节点记录
转换为一个字符串：用连字的组成部分替换连字，并忽略字距。然后在这个字符串上运行
断字算法，再把断字结果转换回一个 ``水平列表，依次拼回段落流中。
请记住，段落中可能含有未装箱的水平内容，其中已经包含连字和字距，
而其中的单词同样参与断字过程。

那些 \type {char} 节点记录的名字多少有些误导，因为它们其实是特定字体中的字形位置，
并非语言学意义上的 ``字符。\type {char} 节点记录内部根本没有语言信息。
相反，语言信息是通过水平列表中的 \type {language whatsit} 节点传递的。

在 \LUATEX\ 中情况则大不相同。你键入的字符总是被转换为 \nod {glyph} 节点记录，
并带有一个特殊子类型，以标明它们是作为语言学字符使用的。\LUATEX\ 把所需的
语言信息存放在这些记录中，但在创建节点时并不做任何与字体相关的处理。
它只保存当前字体的索引，以及该字体中某个字符的引用。

当需要排版一个段落时，\LUATEX\ 先把所有断字点直接插入整个节点列表。
接着处理整个列表中的全部字体信息（创建连字并调整字距），最后调整所有子类型标识，
使这些记录从此成为 ``字形节点。

\section{字符、字形与断字}\label{charsandglyphs}

\topicindex {characters}
\topicindex {glyphs}
\topicindex {hyphenation}

\TEX82（包括 \PDFTEX）区分 \type {char} 节点和 \type
{lig} 节点。前者是简单条目，只含一个 ``字符 字段和一个 ``字体 字段，
与记号存在同一块内存中。后者还包含一个组成部分列表，以及一个子类型，
用以标明该连字是否由词界产生；它与盒子、字距、粘胶等其他节点存放在同一处。

在 \LUATEX\ 中，这两种类型被合并为一个稍大的结构，称为 \nod {glyph} 节点。
除了原有的字符、字体和组成部分字段外，还有另外几个字段，例如我们将在
\ref{glyphnodes} 中看到的 ``attr；这些节点还含有一个子类型，
它编码了四种主要类型和两种附加的“幽灵”类型。对于连字，可以同时置多位
（当一个词只有单个字形时）。

\begin{itemize}\item \type {character}，用于需要断字的字符：最低位（bit 0）置 1。
    
    \item \nod {glyph}，用于特定字体字形：最低位（bit 0）不置位。
    
    \item \type {ligature}，用于构造出的连字，bit 1 置位。
    
    \item \type {ghost}，用于所谓的 ``幽灵对象，bit 2 置位。
    
    \item \type {left}，用于由左词界产生的连字，以及由 \lpr {leftghost}
        产生的幽灵，bit 3 置位。
    
    \item \type {right}，用于由右词界产生的连字，以及由 \lpr {rightghost}
        产生的幽灵，bit 4 置位。
    
\end{itemize}

\nod {glyph} 节点还含有语言数据，分为四项，都是节点创建时的当前值：
\prm {setlanguage}（15~位）、\prm {lefthyphenmin}（8~位）、
\prm {righthyphenmin}（8~位）和 \prm {uchyph}（1~位）。

顺带一提，\LUATEX\ 允许 16383 种不同语言，单词最长可达 256 个字符。
语言信息随每个字符一同保存。你可以把 \prm {firstvalidlanguage}
设为例如~1，从而使语言~0 成为一个被忽略的断字语言。

新原语 \lpr {hyphenationmin} 可用于指定一个词的最小长度。
该值随（当前）语言一同保存。

由于 \prm {uchyph} 的值保存在实际节点中，它的处理与 \TEX82\ 略有不同：
对 \prm {uchyph} 的改动立即生效，而不是在当前段落末尾才生效。

已排版的盒子现在总是把语言信息嵌入节点自身，因此不再可能依赖周围的语言设置。
在 \TEX82\ 中，段落中部出现的 \type {\unhbox0} 之类语句会使用当前段落语言处理盒子，
除非盒子内部发出过 \prm {setlanguage}。在 \LUATEX\ 中，所有语言变量都已冻结。

在传统 \TEX\ 中，断字过程由 \type {lccode} 驱动。在 \LUATEX\ 中我们弱化了这种依赖。
可以有几种策略。什么都不做时，加载模式、设置例外或对列表断字都会使用当前生效的
\type {lccode}。

当你把 \prm {savinghyphcodes} 设为大于零的值时，当前这组 \type {lccode}
会随语言一同保存。这种情况下，之后改动某个 \type {lccode} 不再起作用。
不过你可以用如下方式调整这组代码：

\begin{typing}
\hjcode`a=`a
\end{typing}

这个改动是全局性的——如果你记住断字发生的时刻通常是构造段落或水平盒子之时，
这就说得通了。语言初始化时若 \prm {savinghyphcodes} 为零，则你一开始什么都没有；
否则你已经有了一组代码。

当某个 \lpr {hjcode} 大于 0 但小于 32 时，它表示要使用的长度。在下面的例子中，
我们把一个字符（\type {x}）映射到模式中的另一个字符，并告诉引擎 \type {œ}
算作一个字符。由于传统上 0 本身被保留用于禁止断字，值 32 被当作 0 处理。

下面是一些例子（我们假设使用法语模式）：

\begin{longtable}{l l l}
\toprule & \type{foobar} & \type{foo-bar} \\
\midrule
\type{\hjcode`x=`o}              & \type{fxxbar} & \type{fxx-bar} \\
\type{\lefthyphenmin3}           & \type{œdipus} & \type{œdi-pus} \\
\type{\lefthyphenmin4}           & \type{œdipus} & \type{œdipus}  \\
\type{\hjcode`œ=2}               & \type{œdipus} & \type{œdi-pus} \\
\type{\hjcode`i=32 \hjcode`d=32} & \type{œdipus} & \type{œdipus}  \\
\bottomrule
\end{longtable}

让每个字形都携带这全部信息会带来过大的开销，也会使设置这些代码的过程更加复杂。
我们曾考虑过用 \type {hjcode} 集合的方案，但最终放弃了，因为实践中当前的做法已经足够，
而且无论如何都无法做到兼容。

请注意：这些值总是会被保存到格式文件中，与转储格式时 \prm {savinghyphcodes}
的设置无关。

边界节点通常会标记一个词的结束，这会干扰例如断字插入。为此你可以把
\prm {wordboundary} 用作触发器。下面是几个用法示例：

\begin{typing}
    discrete---discrete
\end{typing}

%（以下为 ConTeXt 演示示例图，已省略。）

在传统 \TEX\ 中，连字构造与断字和换行机制交织在一起。在 \LUATEX\ 中这些阶段是隔离的。
因此我们对（一连串）显式连字符的处理方式有所不同。我们已经添加了一些对断字各方面的
控制，而又一项新的控制涉及自动连字符（例如输入中的 \type {-} 字符）。

当 \lpr {automatichyphenmode} 取值为 0 时，一个连字符会被转换为自动断字符。
它前后的片段不会被断字。一个副作用是，行首连字符可能导致断行，但人们很少遇到这种情况。
设置前、后字符会使这一现象更明显。取值为 1 会避免这个副作用；取值为 2 则不把连字符
转换为断字符。其他选项的实验（例如允许对两侧单词断字）已被放弃。

\begin{typing}
% before-after \par
% before--after \par
% before---after \par
\end{typing}

\begin{typing}
% -before \par
% after- \par
% --before \par
% after-- \par
% ---before \par
% after--- \par
\end{typing}

\begin{typing}
% before-after \par
% before--after \par
% before---after \par
\end{typing}

\begin{typing}
% ConTeXt \framed 演示块，已省略
\end{typing}

% ConTeXt 自动断字模式演示图（automatichyphenmode:1, automatichyphenmode:2），已省略

在下面的讨论中，
我们展示了三个样本的情况：

输入 A、B、C 的演示输出已省略（原 ConTeXt 图）。

与其他单字符命令的原语同伴一样，\prm {-} 命令有一个更冗长的原语版本
\lpr {explicitdiscretionary}；而断字器中通常被截获的字符 \type {-}
（或任何被配置的字符）则可作为 \lpr {automaticdiscretionary} 使用。

\section{主控制循环}

\topicindex {main loop}
\topicindex {hyphenation}

在 \LUATEX\ 的主循环中，几乎所有待排版的输入字符都会被转换为子类型为
``character 的 \nod {glyph} 节点记录，但有少数例外。

\begin{itemize}

\item \prm {accent} 原语创建的节点子类型是 ``glyph 而不是 ``character：
    一个用于实际的重音符号，一个用于被加重音的字符。这样做的首要原因是，
    \TEX82 中的 \prm {accent} 明确依赖于当前字体编码，因此给该原语的名字赋予新含义
    没什么意义，那会使许多旧文档和宏包失效。次要原因是，在 \TEX82\ 中
    \prm {accent} 会禁止对当前单词断字。由于在 \LUATEX\ 中只对 ``character
    节点进行断字，所以可以达到同样的效果。当然，现代具备 \UNICODE\ 感知能力的宏包
    根本不会使用 \prm {accent} 原语，而是直接映射到组合字符。

    \prm {char} 确实发生了这种含义变化：它现在生成带子类型 character 的
    ``glyph 节点。在传统 \TEX\ 中，8 位输入编码、断字与从字体中取字形之间
    存在紧密关系。在 \LUATEX\ 中我们使用 \UTF\ 输入，大多数情况下它直接映射到字体中的
    一个字符，字体引擎中的字形替换除外。如果你想直接访问字体中的任意字形，
    随时可以用 \LUA\ 来做，因为字体是以 \LUA\ 表的形式可用的。

\item 数学模式中处理的所有结果，最终都成为子类型为 ``glyph 的节点。
    事实上，数学处理的结果就是一个由字形、字距、粘胶、罚分、盒子等组成的普通列表。

\item 由 \ALEPH\ 派生的命令 \lpr {leftghost} 和 \lpr {rightghost} 创建第三种
    子类型的节点：``ghost。这些节点在后续所有处理中被完全忽略，
    直到添加字形间字距的阶段才会被处理。

\item 自动断字符的处理方式不同。\TEX82\ 在感知到一个与当前字体 \prm {hyphenchar}
    匹配的输入字符后，会插入一个空断字符。我们认为这种检测方式是错误的：
    是否断字不应依赖于当前字体，它是一种语言属性。
    \footnote {\TEX\ 刚出现时我们还没有 \UNICODE，而局限于 8 位意味着有时不得不在
    支持字符输入、字形渲染和断字之间做出妥协。}

    在 \LUATEX\ 中，工作方式如下：当 \LUATEX\ 感知到一串输入字符与新整数参数
    \prm {exhyphenchar} 的值匹配时，会在这串节点之后插入一个显式断字符。
    初始时 \TEX\ 把 \type {\exhyphenchar=`\-}。顺带一提，这是一个全局参数而非
    语言相关参数，因为根据文档结构而不是文本语言来改变它的值可能更有用。

    在一连串显式连字符之后插入断字符的工作，与其他断字处理同时发生，
    而{\it 不是}在主控制循环内部。

    \LUATEX\ 对 \prm {hyphenchar} 的唯一用途，是检查一个词是否根本应该被纳入断字。
    如果一个词中第一个字符节点所附字体的 \prm {hyphenchar} 为负，则立即放弃对该词的
    断字。这种行为仅是为了向后兼容而添加的，在新的 \LUATEX\ 文档中不应再使用
    \type {\hyphenchar=-1} 作为阻止断字的手段。

\item \prm {setlanguage} 命令不再创建 whatsit。\prm {setlanguage} 的含义被改为一个
    与其他参数一样的整数参数。这个整数参数在创建 \type {\glyph_node} 时被用来
    向字形节点添加语言信息。相应地，\prm {language} 原语也被扩展，使其总是同时更新
    \prm {setlanguage} 的值。

\item \prm {noboundary} 命令（用于在通常会发生词界处理的地方禁止该处理）现在确实会
    创建节点。这些节点是必需的，因为 \prm {noboundary} 命令在输入流中的确切位置
    必须保留到连字和字体处理阶段之后。

\item 代码中不再有 \type {main_loop} 标号。请记住，\TEX82\ 在向水平列表添加
    \type {char_nodes} 时做了大量处理。出于速度考虑，它把那段处理代码放在
    ``主控制 循环之外，只有任何 ``词 的第一个字符才由那个
    ``主控制 循环处理。在 \LUATEX\ 中不再需要这样做（所有重活都在后面完成），
    那些（现在非常短小的）字符处理代码已经移回内联。当打开 \prm {tracingcommands} 时，
    这一点可以看出来：报告的是整个词，而不仅仅是首字符。

\end{itemize}

由于我们倾向于把硬编码的行为改为可配置，又添加了几个新原语：

\begin{typing}
\hyphenpenaltymode
\automatichyphenpenalty
\explicithyphenpenalty
\end{typing}

第一个参数对自动断字符（即由 \prm {exhyphenchar} 产生的断字符）有以下影响：

\begin{longtable}{c l l}
\toprule 模式     & 自动断字符 \type {-}      & 显式断字符 \prm{-}         \\
\midrule
\type{0} & \prm {exhyphenpenalty}        & \prm {exhyphenpenalty}        \\
\type{1} & \prm {hyphenpenalty}          & \prm {hyphenpenalty}          \\
\type{2} & \prm {exhyphenpenalty}        & \prm {hyphenpenalty}          \\
\type{3} & \prm {hyphenpenalty}          & \prm {exhyphenpenalty}        \\
\type{4} & \lpr {automatichyphenpenalty} & \lpr {explicithyphenpenalty}  \\
\type{5} & \prm {exhyphenpenalty}        & \lpr {explicithyphenpenalty}  \\
\type{6} & \prm {hyphenpenalty}          & \lpr {explicithyphenpenalty}  \\
\type{7} & \lpr {automatichyphenpenalty} & \prm {exhyphenpenalty}        \\
\type{8} & \lpr {automatichyphenpenalty} & \prm {hyphenpenalty}          \\
\bottomrule
\end{longtable}

其他值则执行我们在 \LUATEX\ 中一直以来的做法：插入 \prm {exhyphenpenalty}。

\section{加载断字模式与例外词}\label{patternsexceptions}

\topicindex {hyphenation}
\topicindex {hyphenation+patterns}
\topicindex {hyphenation+exceptions}
\topicindex {patterns}
\topicindex {exceptions}

尽管我们保留了传统的断字方法（它依然更优越），但 \LUATEX\ 中断字算法的实现与
\TEX82\ 大不相同。

\prm {patterns} 的参数经过展开后必须是合法的 \UTF8，各模式之间以空格分隔，
不允许出现 \prm {char} 或 \prm {chardef} 命令。当前实现相当严格，会拒绝所有
非-\UNICODE\ 字符。同样，\prm {hyphenation} 展开后的参数也必须是合法的 \UTF8，
但这里额外提供了一点语法：

\begin{itemize}
\item 花括号中的三组参数（\type {{}{}{}}）表示所需的复合断字符，
    参数与普通文档输入中 \prm {discretionary} 命令的参数相同。

\item 一个 \type {-} 表示所需的简单断字符，参见普通文档输入中的
    \type {\-} 和 \type {\discretionary{-}{}{}}。

\item 内部命令名会被忽略。这条规则专门为 \prm {discretionary} 提供，
    但它也有助于处理可能混入的 \prm {relax} 命令。

\item 一个 \type {=} 表示文档输入中的一个（非断字符）连字符。

\end{itemize}

展开后的参数首先被转换回一个空格分隔的字符串，同时丢弃内部命令名。
然后由一个例程把这个字符串转换为字典：通过转换其他列出的项来创建键|-|值对。
需要重要注意的是，例外字典中的键总是可以由值推导出来。下面是几个例子：

\begin{longtable}{l l l}
\toprule 取值                  & 推导键（输入） & 效果 \\
\midrule
\type {ta-ble}         & table               & \type {ta\-ble} ($=$ \type {ta\discretionary{-}{}{}ble}) \\
\type {ba{k-}{}{c}ken} & backen              & \type {ba\discretionary{k-}{}{c}ken} \\
\bottomrule
\end{longtable}

得到的模式和例外字典会保存在 \prm {language} 当前值所对应的语言代码之下。

在表格的最后一行中，你可以看到值中没有 \prm {discretionary} 命令：
在基于 \TEX\ 的输入语法中，该命令是可选的。其背后的原因是，可以设想把一整本
单词字典存为纯文本文件，再用 \LUA\ \type {lang} 库中的某个函数加载进 \LUATEX。
这种加载方式比走 \TEX\ 语言原语快不少，但如果在加载时还要解释命令序列，
那部分（大部分）速度优势就会丧失。

可以用 \type {{-}{}{-}} 为复合词中的显式连字符指定额外的断字点
（如有需要，把 \type {-} 换成实际的显式连字符）。例如，下面这条与单词
``multi|-|word|-|boundaries 匹配，并允许在 ``boun 和
``daries 之间额外断行：

\begin{typing}
{}{-}word{-}{}{-}boun-daries}
\end{typing}

\ETEX\ 扩展 \prm {savinghyphcodes} 的动机，是断字过去严重依赖于字体编码。
在 \LUATEX\ 中情况已不再如此，对应的原语基本被忽略了。因为我们现在有了
\lpr {hjcode}，大小写相关代码可以专门用于 \prm {uppercase} 和 \prm {lowercase}。

例外词中的三对花括号模式可能有些出人意料，所以我们通过例子来解释。
模式 \type {foo{}{}{x}bar} 会创建查找键 \type {fooxbar}，而模式
\type {foo{}{}{}bar} 创建 \type {foobar}。然后，当命中发生时，
有一个替换文本（\type {x}）或没有。由于我们在断字符节点中引入了罚分，
例外语法现在还可以带罚分说明。方括号中的值是 \lpr {exceptionpenalty} 的倍数。
这里我们把它设为 10000，因此在例子中实际得到 30000。

% ConTeXt \ShowSample 演示宏（使用 \bTABLE/\bTR/\bTD 等 ConTeXt 表格命令），已省略。
% 以下为原示例断字串：
% x{a-}{-b}{}x{a-}{-b}{}x{a-}{-b}{}x{a-}{-b}{}xx
% x{a-}{-b}{}x{a-}{-b}{}[3]x{a-}{-b}{}[1]x{a-}{-b}{}xx
% z{a-}{-b}{z}{a-}{-b}{z}{a-}{-b}{z}{a-}{-b}{z}z
% z{a-}{-b}{z}{a-}{-b}{z}[3]{a-}{-b}{z}[1]{a-}{-b}{z}z

\section{应用断字}

\topicindex {hyphenation+how it works}
\topicindex {hyphenation+discretionaries}
\topicindex {discretionaries}

\LUATEX\ 在单词中插入断字符所使用的内部结构与 \TEX82\ 非常不同，这意味着在处理上
也有一些值得注意的差异。

首要一点是，断字过程中不涉及 ``压缩 trie。算法仍然读取由 \PATGEN\ 生成的
模式文件，但 \LUATEX\ 使用有限状态哈希把模式与待断字单词进行匹配。该算法基于
\OPENOFFICE\ 所用的 ``libhnj 库，而后者又受 \TEX\ 启发。

由于实现方式的原因，\LUATEX\ 与 \TEX82\ 之间存在一些差异：

\begin{itemize}\item \LUATEX\ 可以愉快地对整个 \UNICODE\ 字符范围断字。

\item 模式和例外字典的大小只受可用内存限制，所有分配都是动态进行的。
    \type {texmf.cnf} 中与 trie 相关的设置会被忽略。

\item 由于不存在 ``trie 准备 阶段，语言模式永远不会被冻结。这意味着原语
    \prm {patterns}（及其 \LUA\ 对应项 \type {lang.patterns}）可以在任何时候使用，
    而不仅仅在 ini\TEX\ 中。

\item 格式文件中只保存 \prm {patterns} 和 \prm {hyphenation} 的字符串表示。
    在加载格式时，它们只是被重新求值。由此可见，没有真正理由在格式文件中预加载语言。
    事实上，这样做通常并不好。更聪明的做法是，直到真正需要某种模式时才加载它。

\item \LUATEX\ 在创建隐含断字符时使用语言相关变量 \lpr {prehyphenchar} 和
    \lpr {posthyphenchar}，而不是 \TEX82 的 \prm {hyphenchar}；
    并使用语言相关变量 \lpr {preexhyphenchar} 和 \lpr {postexhyphenchar} 的值
    来创建显式断字符（而不是 \TEX82 的空断字符）。

\item 与断字相关的两个计数器 \prm {hyphenpenalty} 和 \prm {exhyphenpenalty} 的值，
    现在保存在断字符节点中。这允许对显式 \prm {discretionary} 命令进行局部覆盖。
    应用断字遍历时所使用的是当时的当前值。不使用回调时，这与传统 \TEX\ 兼容。
    当你调用 \LUA\ 函数 \type {lang.hyphenate} 时，使用的也是当前值。

\item 断字例外字典以键|-|值哈希的形式维护，而且同样是动态的，
    因此 \type {hyph_size} 设置也不再被使用。

\end{itemize}

因为我们在断字符节点中保存了罚分，\prm {discretionary} 命令被扩展为可以接受一个
可选的罚分说明，于是你可以这样做：

\begin{typing}
\hsize1mm
1:foo{\hyphenpenalty 10000\discretionary{}{}{}}bar\par
2:foo\discretionary penalty 10000 {}{}{}bar\par
3:foo\discretionary{}{}{}bar\par
\end{typing}

结果如下（演示输出已省略）。

\libindex {prehyphenchar}
\libindex {posthyphenchar}
\libindex {preexhyphenchar}
\libindex {postexhyphenchar}

\begin{typing}
<number> n = lang.prehyphenchar(<language> l)
lang.prehyphenchar(<language> l, <number> n)

<number> n = lang.posthyphenchar(<language> l)
lang.posthyphenchar(<language> l, <number> n)
\end{typing}

这两个函数用于获取或设置该语言中隐含断字的 ``断前 和 ``断后
连字符。初始值为十进制 45（连字符）和十进制~0（表示为空）。

\begin{typing}
<number> n = lang.preexhyphenchar(<language> l)
lang.preexhyphenchar(<language> l, <number> n)

<number> n = lang.postexhyphenchar(<language> l)
lang.postexhyphenchar(<language> l, <number> n)
\end{typing}

这两个函数用于获取或设置该语言中显式断字的 ``断前 和 ``断后
连字符。两者初始值均为十进制~0（表示为空）。

\subsection {\type {hyphenate}}

\libindex {hyphenate}

下一个调用在节点列表中插入断字点（断字符节点）。如果以 \type {tail} 为参数，
则处理到该节点即停止。目前，如果 \type {head}（以及指定的 \type {tail}）
是合法节点，无论是否可能有其他错误，\type {success} 总是为真。

\begin{typing}
<boolean> success = lang.hyphenate(<node> head)
<boolean> success = lang.hyphenate(<node> head, <node> tail)
\end{typing}

断字只作用于 ``character，即所有字形节点中子类型取值为 \type {1} 的那个
特殊子类型。子类型不同的字形模式不被处理。更多细节见
\ref{charsandglyphs}。

\subsection {\type {[set|get]hjcode}}

\libindex {sethjcode}
\libindex {gethjcode}

以下两条命令可用于设置或查询 hj code：

\begin{typing}
lang.sethjcode(<language> l, <number> char, <number> usedchar)
<number> usedchar = lang.gethjcode(<language> l, <number> char)
\end{typing}

当你设置一个 hjcode 时，当前这组代码会被初始化，除非该组已经因为
\prm {savinghyphcodes} 大于零而初始化过。

% \parindent0pt \hsize=1.1cm
% 12-34-56 \par
% 12-34-\hbox{56} \par
% 12-34-\vrule width 1em height 1.5ex \par
% 12-\hbox{34}-56 \par
% 12-\vrule width 1em height 1.5ex-56 \par
% \hjcode`\1=`\1 \hjcode`\2=`\2 \hjcode`\3=`\3 \hjcode`\4=`\4 \vskip.5cm
% 12-34-56 \par
% 12-34-\hbox{56} \par
% 12-34-\vrule width 1em height 1.5ex \par
% 12-\hbox{34}-56 \par
% 12-\vrule width 1em height 1.5ex-56 \par
